\section{Related Work}

\subsection{Affect as a geometric target}

The use of emotion as a probe of representation geometry has a useful external
anchor.  Psychological accounts often organize affect continuously by valence
(pleasantness) and arousal (activation), rather than treating every emotion as
an unrelated category \cite{russell_circumplex_1980,posner_circumplex_2005}.
This does not imply that a language model implements human affective
mechanisms.  It does, however, provide a testable expectation: if a model has
coherent representations of emotion concepts, related concepts may cluster and
their leading dimensions may align with these established descriptors.
Recent analyses report such structure in LLM hidden states.  Choi and Weber
\cite{choi_latent_2026} find affective representations that align with
valence--arousal models and are approximately linearly recoverable, whereas
Reichman et al.\ \cite{reichman_emotions_2026} report a low-dimensional,
cross-lingual emotional subspace.  Sun et al.\ \cite{sun_valence-arousal_2026}
further report a circular valence--arousal geometry and monotonic affective
control in several open-weight models.  Our PCA and cosine-similarity analyses
are qualitative tests in this tradition; they do not estimate an external
valence--arousal correlation and therefore should not be interpreted as a
quantitative validation of either model.

\subsection{Activation steering and interpretation}

Activation steering modifies an internal activation by adding a direction
associated with a desired property, then measures the resulting change in
output.  It is a causal intervention on the model computation, but a successful
intervention alone does not establish that the direction is a unique or
complete representation of the named concept.  This distinction matters
because steering vectors can fall outside the distribution on which common
interpretability tools are trained and can have meaningful negative feature
components \cite{mayne_can_2024}.  The residual-stream setting used here is
also compatible with work showing that Transformer components combine
interpretable patterns through residual connections before producing output
distributions \cite{geva_transformer_2021}.  We use Logit Lens projections and
steered generations as complementary diagnostic evidence, rather than treating
either one as a semantic ground truth.

\subsection{Emotion vectors in language models}

The direct precursor to this study is the analysis of Claude Sonnet 4.5 by
Sofroniew et al.\ \cite{sofroniew_emotion_2026}.  They extracted representations
for 171 emotion concepts, showed that the representations tracked contextual
emotion and exhibited valence--arousal structure, and tested their functional
effects with interventions.  Their term \emph{functional emotions} refers to
representations that influence behavior; it explicitly does not entail a claim
of subjective experience.

The question of transfer beyond Claude is now beginning to receive direct
attention.  Hsieh \cite{hsieh_replicating_2026} released an early community
replication on Gemma 4 E4B.  In a systematic open-weight study, van der Ben et
al.\ \cite{ben_where_2026} recover valence geometry in Gemma 4 E4B-it and
Apertus-8B-Instruct-2509, while finding that the layer at which valence emerges
and the apparent arousal structure depend on the model and extraction corpus.
These results make model choice, layer selection, and corpus construction
central experimental variables rather than incidental implementation details.
Our study complements this evidence by applying a fixed, inexpensive pipeline
to GPT-2 Medium and Gemma 4 E2B, and by contrasting a nine-category set with a
20-category set.  It is therefore best understood as a partial replication and
an exploratory robustness check, not as a test of the full behavioral and
safety claims in the original work.
